%% CV for Norlys - Cloud Engineer (Dansk)
%% Compiled with: lualatex

\documentclass[11pt,a4paper,sans]{moderncv}
\moderncvstyle{banking}
\moderncvcolor{blue}

\renewcommand*{\firstnamestyle}[1]{{\fontsize{34}{36}\bfseries\upshape\color{color1}#1}}
\renewcommand*{\lastnamestyle}[1]{{\fontsize{34}{36}\bfseries\upshape\color{color1}#1}}
\renewcommand*{\sectionstyle}[1]{{\sectionfont\color{color1}#1}}

\usepackage[utf8]{inputenc}
\usepackage{hyperref}
\hypersetup{
    colorlinks=true,
    linkcolor=blue,
    filecolor=magenta,
    urlcolor=blue,
    pdftitle={Jacob El-Omar - CV},
    pdfpagemode=FullScreen,
}
\usepackage[scale=0.80]{geometry}
\usepackage{import}

\name{Jacob}{El-Omar}
\address{Aalborg, Danmark}{}{}
\phone[mobile]{+45 53 35 27 82}
\email{elomarjc@gmail.com}
\extrainfo{\href{https://linkedin.com/in/jacob-el-omar}{LinkedIn} \textbullet{} \href{https://github.com/elomarjc}{GitHub}}

\begin{document}

\makecvtitle

% ============================================================
%     PROFIL
% ============================================================

\vspace{6pt}
\small{Nyuddannet civilingeniør (cand.polyt.) i elektroniske systemer fra Aalborg Universitet med stærke kompetencer inden for programmering (Python og C\#) og distribuerede logikker. Erfaren i at opbygge backend-API'er, databasemodeller samt automatisere deployment-processer. En nysgerrig og pragmatisk holdspiller, der motiveres af at drifte og videreudvikle platformsinfrastruktur i jeres Cloud Development-team.}

% ============================================================
%     FAGLIGE KOMPETENCER
% ============================================================

\section{Faglige Kompetencer}
\vspace{1pt}
\begin{itemize}

\item \textbf{Programmering \& Scripting}: Python, C\#, .NET Core, JavaScript, TypeScript, Node.js, bash scripting.

\item \textbf{Infrastruktur \& Platformsforståelse}: Erfaring med CI/CD processer (GitHub Actions), REST-API'er, datamodellering, og interesse for Infrastructure-as-Code (Terraform).

\item \textbf{Databaser \& Systemintegration}: SQL (PostgreSQL, SQLite, MySQL, Supabase), hændelsesbaseret software og API-integrationer.

\item \textbf{Overvågning \& Fejlsøgning}: Systematiske testmetoder, datavisualisering, samt logs- og fejlsøgning på servere og databaser.

\item \textbf{Samarbejde \& Læring}: Uddannet efter AAU's problembaserede læringsmodel (PBL), som sikrer stærke formidlings- og teamwork-kompetencer.

\end{itemize}

% ============================================================
%     ERHVERVSERFARING
% ============================================================

\section{Erhvervserfaring}
\vspace{3pt}
\begin{itemize}

\item{\cventry{Jan 2026--Nu}{IT- \& Webudvikler (Frivillig)}{Dawah Danmark}{Aalborg, Danmark}{}{\vspace{1pt}
\begin{itemize}
    \item Vedligeholdt og driftede serverplatformen for organisationens hjemmeside.
    \item Konfigurerede plug-ins, overvågede systemstabilitet og løste ydeevneproblemer.
\end{itemize}}}

\vspace{3pt}

\item{\cventry{Sep 2024--Jun 2025}{AI/ML Engineer (Praktik)}{Ai Health Highway}{Remote}{}{\vspace{1pt}
\begin{itemize}
    \item Udviklede og testede realtids-signalfiltrering i Python til medicinske IoT-enheder.
    \item Samarbejdede om hardware-software integrationskrav med globale forskerhold.
    \item Udarbejdede teknisk dokumentation og testcases for at sikre softwarens overholdelse af standarder.
\end{itemize}}}

\vspace{3pt}

\item{\cventry{Jun 2023--Aug 2024}{Lager- og Logistikmedarbejder}{Lyn Express}{Aalborg, Danmark}{}{\vspace{1pt}
\begin{itemize}
    \item Koordinerede varemodtagelse, sortering og databaseopdateringer under højt tidspres.
\end{itemize}}}

\end{itemize}

% ============================================================
%     UDDANNELSE
% ============================================================

\section{Uddannelse}
\vspace{1pt}
\begin{itemize}

\item{\cventry{Sep 2023--Jun 2025}{Civilingeniør, cand.polyt. (Elektroniske Systemer)}{Aalborg Universitet}{Aalborg, Denmark}{}{\vspace{1pt}
Speciale: ``Adaptive Noise Cancellation For Electronic Stethoscopes.'' Fokuseret på støjfiltrering, signalbehandling og distribueret software.
}}

\vspace{3pt}

\item{\cventry{Sep 2020--Aug 2023}{Bachelor i teknisk videnskab (Elektronik og IT)}{Aalborg University}{Aalborg, Danmark}{}{\vspace{1pt}
Specialisering i reguleringsteknik. Projekter indeholdt automatisk faldsikringsprogrammering i C på Cypress PSoC og regulering til robotsystemer.
}}

\end{itemize}

% ============================================================
%     PROJEKTER (UDVALGTE)
% ============================================================

\section{Udvalgte Projekter}
\vspace{1pt}
\begin{itemize}
\item{\textbf{Grow a Fish}: Udvikling af et fuldt spil i Flutter/Dart. Oprettede en Supabase PostgreSQL-database til håndtering af multiplayer-data via REST-API.}
\item{\textbf{AAUSAT6 ADCS Design}: Udviklede simulations- og styringsmoduler for satellitstabilisering i MATLAB/Simulink.}
\end{itemize}

% ============================================================
%     SPROG
% ============================================================

\section{Sprog}
\vspace{1pt}
\begin{itemize}
\item Dansk (modersmål), Engelsk (flydende).
\end{itemize}

% ============================================================
%     REFERENCER
% ============================================================

\section{References}
\vspace{1pt}
\begin{itemize}
\item{Afleveres efter aftale.}
\end{itemize}

\end{document}
